\documentclass[11pt,a4paper]{article}
% preamble.tex — estilo común de todos los apuntes Froth.
% Cada apunte hace % preamble.tex — estilo común de todos los apuntes Froth.
% Cada apunte hace % preamble.tex — estilo común de todos los apuntes Froth.
% Cada apunte hace \input{preamble} justo después de \documentclass.
\usepackage[margin=2.4cm]{geometry}
\usepackage{xcolor}
\definecolor{litblue}{RGB}{31,79,121}
\definecolor{litgreen}{RGB}{27,94,32}
\definecolor{litorange}{RGB}{230,81,0}
\usepackage{parskip}
\usepackage{titlesec}
\titleformat{\section}{\Large\bfseries\color{litblue}}{\thesection}{0.6em}{}
\titleformat{\subsection}{\large\bfseries\color{litblue}}{\thesubsection}{0.6em}{}
\usepackage{tcolorbox}
\usepackage{fancyhdr}
\pagestyle{fancy}
\fancyhf{}
\fancyhead[L]{\small\color{litblue}Froth — Apuntes de ML}
\fancyhead[R]{\small\thepage}
\renewcommand{\headrulewidth}{0.4pt}
\usepackage[colorlinks=true,linkcolor=litblue,urlcolor=litblue]{hyperref}

% Cabecera de cada apunte: \cabecera{numero}{titulo}
\newcommand{\cabecera}[2]{%
  \begin{tcolorbox}[colback=litblue,coltext=white,colframe=litblue,arc=2mm]
    {\Large\bfseries Apunte #1 — #2}\\[3pt]
    {\small Proyecto Froth · aprender Machine Learning construyendo}
  \end{tcolorbox}\vspace{0.8em}}

% Cajas temáticas
\newtcolorbox{concepto}[1]{colback=blue!4,colframe=litblue,title={Concepto: #1},arc=1.5mm}
\newtcolorbox{practica}{colback=green!5,colframe=litgreen,title={Pruébalo tú (teclado en mano)},arc=1.5mm}
\newtcolorbox{ojo}{colback=orange!8,colframe=litorange,title={Ojo / error típico},arc=1.5mm}
\newtcolorbox{resumen}{colback=gray!8,colframe=gray!60,title={En una frase},arc=1.5mm}

\newcommand{\codigo}[1]{\texttt{#1}}
 justo después de \documentclass.
\usepackage[margin=2.4cm]{geometry}
\usepackage{xcolor}
\definecolor{litblue}{RGB}{31,79,121}
\definecolor{litgreen}{RGB}{27,94,32}
\definecolor{litorange}{RGB}{230,81,0}
\usepackage{parskip}
\usepackage{titlesec}
\titleformat{\section}{\Large\bfseries\color{litblue}}{\thesection}{0.6em}{}
\titleformat{\subsection}{\large\bfseries\color{litblue}}{\thesubsection}{0.6em}{}
\usepackage{tcolorbox}
\usepackage{fancyhdr}
\pagestyle{fancy}
\fancyhf{}
\fancyhead[L]{\small\color{litblue}Froth — Apuntes de ML}
\fancyhead[R]{\small\thepage}
\renewcommand{\headrulewidth}{0.4pt}
\usepackage[colorlinks=true,linkcolor=litblue,urlcolor=litblue]{hyperref}

% Cabecera de cada apunte: \cabecera{numero}{titulo}
\newcommand{\cabecera}[2]{%
  \begin{tcolorbox}[colback=litblue,coltext=white,colframe=litblue,arc=2mm]
    {\Large\bfseries Apunte #1 — #2}\\[3pt]
    {\small Proyecto Froth · aprender Machine Learning construyendo}
  \end{tcolorbox}\vspace{0.8em}}

% Cajas temáticas
\newtcolorbox{concepto}[1]{colback=blue!4,colframe=litblue,title={Concepto: #1},arc=1.5mm}
\newtcolorbox{practica}{colback=green!5,colframe=litgreen,title={Pruébalo tú (teclado en mano)},arc=1.5mm}
\newtcolorbox{ojo}{colback=orange!8,colframe=litorange,title={Ojo / error típico},arc=1.5mm}
\newtcolorbox{resumen}{colback=gray!8,colframe=gray!60,title={En una frase},arc=1.5mm}

\newcommand{\codigo}[1]{\texttt{#1}}
 justo después de \documentclass.
\usepackage[margin=2.4cm]{geometry}
\usepackage{xcolor}
\definecolor{litblue}{RGB}{31,79,121}
\definecolor{litgreen}{RGB}{27,94,32}
\definecolor{litorange}{RGB}{230,81,0}
\usepackage{parskip}
\usepackage{titlesec}
\titleformat{\section}{\Large\bfseries\color{litblue}}{\thesection}{0.6em}{}
\titleformat{\subsection}{\large\bfseries\color{litblue}}{\thesubsection}{0.6em}{}
\usepackage{tcolorbox}
\usepackage{fancyhdr}
\pagestyle{fancy}
\fancyhf{}
\fancyhead[L]{\small\color{litblue}Froth — Apuntes de ML}
\fancyhead[R]{\small\thepage}
\renewcommand{\headrulewidth}{0.4pt}
\usepackage[colorlinks=true,linkcolor=litblue,urlcolor=litblue]{hyperref}

% Cabecera de cada apunte: \cabecera{numero}{titulo}
\newcommand{\cabecera}[2]{%
  \begin{tcolorbox}[colback=litblue,coltext=white,colframe=litblue,arc=2mm]
    {\Large\bfseries Apunte #1 — #2}\\[3pt]
    {\small Proyecto Froth · aprender Machine Learning construyendo}
  \end{tcolorbox}\vspace{0.8em}}

% Cajas temáticas
\newtcolorbox{concepto}[1]{colback=blue!4,colframe=litblue,title={Concepto: #1},arc=1.5mm}
\newtcolorbox{practica}{colback=green!5,colframe=litgreen,title={Pruébalo tú (teclado en mano)},arc=1.5mm}
\newtcolorbox{ojo}{colback=orange!8,colframe=litorange,title={Ojo / error típico},arc=1.5mm}
\newtcolorbox{resumen}{colback=gray!8,colframe=gray!60,title={En una frase},arc=1.5mm}

\newcommand{\codigo}[1]{\texttt{#1}}

\begin{document}
\cabecera{44}{Tripletes de citación: entrenar como SPECTER --- y el veredicto invertido}

\section{Glosario en palabras simples}

\begin{itemize}
  \item \textbf{Fine-tuning}: tomar un modelo ya entrenado (SPECTER) y darle un
        ``curso de especialización'' con datos de TU dominio.
  \item \textbf{Tripleta}: un ejemplo de entrenamiento con 3 piezas: \emph{ancla} (un
        paper), \emph{positivo} (uno que debe quedar CERCA) y \emph{negativo} (uno que
        debe quedar LEJOS).
  \item \textbf{Hard negative}: un negativo DIFÍCIL --- se parece al ancla pero no es lo
        que buscas. Enseña mucho más que un negativo obvio.
  \item \textbf{Época}: una pasada completa por todos los ejemplos de entrenamiento.
  \item \textbf{DBCV}: nota de calidad de un clustering por densidad ($-1$ a $1$):
        premia grupos densos por dentro y bien separados entre sí.
  \item \textbf{Doble ciego}: evaluar sin saber qué opción es cuál, para que el juicio
        no se contamine.
\end{itemize}

\section{La pregunta que hiciste}

\emph{``¿Deberíamos poder mejorar el fine-tune, o SPECTER es perfecto?''} La respuesta
corta: SPECTER no es perfecto, pero nuestros fine-tunes anteriores usaban un método
débil. v1 y v2 entrenaban con pares (título $\leftrightarrow$ abstract): eso solo enseña
``los títulos van con sus abstracts'', y arrastra todo al centro de masa del corpus ---
por eso v2 (3$\times$ datos) NO superó a v1 (apunte 41). El SPECTER original se entrena
distinto: con \textbf{tripletes de citación} --- si el paper A cita al B, B debe quedar
cerca de A; un paper que A no cita debe quedar lejos.

\section{Mini-ejemplo de juguete}

Ancla: ``Flotación de micas con colectores catiónicos''. Positivo: el paper de colectores
que ELLA CITA. Hard negative: un paper citado por el positivo pero NO por el ancla (a dos
saltos: pariente del pariente, pero no tuyo). En cada paso, el entrenamiento MIDE las
distancias y empuja: positivo un poco más cerca, negativo un poco más lejos. Repite
100.000 veces y el espacio entero se reorganiza según quién-cita-a-quién.

La cosecha sin techo hizo posible esto a escala: de 833 tripletes (corpus viejo) a
\textbf{150.885} (100.190 con hard negative) --- 181$\times$ más señal, porque con
190.000 papers muchas referencias caen DENTRO del corpus.

\section{El entrenamiento (y sus perillas medidas)}

Primer intento: 2 épocas $\times$ secuencias de 512 tokens $=$ 17 horas estimadas en tu
RTX 5060 (3,2 s/paso, medido). Decisión tuya: 1 época + 256 tokens (los abstracts
$\sim$200 tokens caben casi enteros) + checkpoints cada 1.000 pasos (lección del apunte
42: una interrupción cuesta minutos, no todo). Resultado: 3,4 pasos/segundo --- terminó
en $\sim$45 minutos. El modelo: \texttt{specter-mineral-v3}.

\section{El veredicto en dos rondas}

\textbf{Ronda 1 --- la métrica (DBCV, 22 corpus limpios, muestra fija de 1.500):}

\begin{center}
\begin{tabular}{lcc}
 & gana a base & delta medio \\
\hline
v1 (pares, chico)   & 14/22 & $+0.013$ \\
v2 (pares, grande)  & 9/22  & $-0.022$ \\
v3 (citación)       & 12/22 & $\mathbf{+0.023}$ \\
v3 vs v2 directo    & 13/22 & $+0.046$ \\
\end{tabular}
\end{center}

Tu pregunta quedó respondida: el MÉTODO correcto gana donde la CANTIDAD falló (v2 pierde
contra base; v3 es el único con mejora media clara). Y un hallazgo extra: en los corpus
sucios viejos v1 ganaba 16/22 con $+0.093$; al limpiar, TODAS las ventajas se encogieron
--- parte de la ``mejora'' anterior era saber agrupar la basura.

\textbf{Ronda 2 --- el experto (tú, doble ciego, 4 mapas anónimos):} rankeaste
B$>$D$>$C$>$A\ldots{} que al revelar la llave resultó ser \textbf{v1 $>$ v2 $>$ base $>$
v3}. ¡El ranking casi ESPEJO de la métrica! El favorito del DBCV (v3, 0.377) fue tu
último lugar.

\begin{concepto}{geometría $\neq$ narrativa}
El DBCV mide si los grupos son densos y separados (geometría). Tú mides si el mapa
CUENTA la historia del campo (narrativa). v3 aprendió de citas, así que agrupa
\emph{comunidades de citación} --- quién cita a quién: escuelas, métodos, revistas ---
que forman blobs preciosos para la métrica pero no siempre coinciden con TEMAS legibles.
v1, con su ajuste suave, conserva la organización semántica que un experto lee.
\end{concepto}

\section{Las decisiones (tuyas, como juez)}

\begin{itemize}
  \item \textbf{Producción $\to$ v1}: lo elegiste a ciegas DOS veces (2026-07-12 y
        2026-07-16) y gana a base 14/22. Tu instalación usa v1 local; la versión pública
        cae automáticamente a SPECTER de fábrica (sin fricción, sin descargas manuales).
  \item \textbf{v3 no se bota --- se re-propósito}: su espacio ``citacional'' es justo la
        señal que necesita el futuro modo semilla (``papers similares a ESTE'', estilo
        ResearchRabbit): ahí las comunidades de citación SON lo que buscas.
\end{itemize}

\begin{ojo}
Si solo hubiéramos mirado la métrica, habríamos puesto en producción el modelo que el
experto lee PEOR. Si solo hubiéramos mirado al experto, no sabríamos POR QUÉ v3 agrupa
raro (ni que es oro para otra tarea). Máquina propone, experto valida, y las dos
miradas juntas valen más que cualquiera sola.
\end{ojo}

\begin{resumen}
La receta real de SPECTER (tripletes de citación + hard negatives) aplicada a TUS
190.000 papers: 150.885 tripletes, 45 min de GPU, v3. La métrica lo corona ($+0.023$
sobre base, único positivo claro; método $>$ cantidad confirmado); tu doble ciego lo
manda de último y corona a v1 --- ranking espejo. Lección doble: (1) el método de
entrenamiento importa más que el volumen de datos; (2) la calidad geométrica no es
calidad narrativa: elige el modelo según la TAREA (v1 para mapas legibles, v3 como
candidato para similitud por citación). ``No mejora para esto'' también es resultado
--- y encontrarle su tarea correcta, también.
\end{resumen}

\end{document}
